
\section{含时微扰理论：系统复习与公式手册}
\label{sec:tdpt-review}

\subsection{第二激发态 $|m|=2$ 的二级修正}

\begin{derivationbox}[第二激发态 $|m|=2$ 的二级修正]
对 $|2\rangle$，非零矩阵元连接到 $|1\rangle$ 和 $|3\rangle$：
\[
E_2^{(2)} = \frac{|\langle 1|H'|2\rangle|^2}{E_2^{(0)} - E_1^{(0)}} + \frac{|\langle 3|H'|2\rangle|^2}{E_2^{(0)} - E_3^{(0)}}.
\]
代入：
\[
E_2^{(2)} = \frac{(erE/2)^2}{\frac{4\hbar^2}{2mr^2} - \frac{\hbar^2}{2mr^2}} + \frac{(erE/2)^2}{\frac{4\hbar^2}{2mr^2} - \frac{9\hbar^2}{2mr^2}} = \frac{mr^4 e^2 E^2}{6\hbar^2} - \frac{mr^4 e^2 E^2}{10\hbar^2} = \frac{mr^4 e^2 E^2}{15\hbar^2}.
\]

\begin{center}
\begin{tabular}{ccccc}
\hline
能级 & 状态 & 简并度 & $E^{(1)}$ & $E^{(2)}$ \\
\hline
基态 ($m=0$) & $|0\rangle$ & 1 & 0 & $-\dfrac{mr^4 e^2 E^2}{\hbar^2}$ \\
第一激发态 ($|m|=1$) & $|\pm 1\rangle$ & 2 & 0 & $\dfrac{mr^4 e^2 E^2}{3\hbar^2}$ \\
第二激发态 ($|m|=2$) & $|\pm 2\rangle$ & 2 & 0 & $\dfrac{mr^4 e^2 E^2}{15\hbar^2}$ \\
\hline
\end{tabular}
\end{center}
\end{derivationbox}

\subsection{基态的二级修正（补充）}

\begin{derivationbox}[基态的二级修正（补充）]
虽然题目只要求激发态，但基态 ($m=0$) 的修正也很有教学意义。$|0\rangle$ 只与 $|\pm 1\rangle$ 耦合：
\[
E_0^{(2)} = \frac{|\langle 1|H'|0\rangle|^2}{E_0^{(0)} - E_1^{(0)}} + \frac{|\langle -1|H'|0\rangle|^2}{E_0^{(0)} - E_{-1}^{(0)}} = 2 \times \frac{(erE/2)^2}{0 - \dfrac{\hbar^2}{2mr^2}} = -\frac{mr^4 e^2 E^2}{\hbar^2}.
\]
负号符合物理直觉：感生电偶极矩与外电场反向，能量降低。
\end{derivationbox}

\subsection{为什么需要含时微扰？}

定态微扰处理的是不随时间变化的势对能级的修正（如 Stark 效应、精细结构）。含时微扰处理的是随时间变化的势导致的量子态之间的跃迁（如原子吸收/发射光子、粒子在时变电磁场中的激发）。

\subsection{核心公式速查卡}

\begin{cheatbox}[核心公式速查卡]
一级近似跃迁概率幅：
\[
a_{fi}(t) = \frac{1}{i\hbar} \int_0^t \langle f|H'(t')|i\rangle e^{i\omega_{fi}t'} \, dt'
\]
其中 $\omega_{fi} = (E_f - E_i)/\hbar$。

跃迁概率：$P_{i\to f}(t) = |a_{fi}(t)|^2$。
\end{cheatbox}

\subsection{相互作用绘景与基本方程}

\begin{derivationbox}[相互作用绘景与基本方程]
薛定谔绘景中 $H = H_0 + H'(t)$。引入相互作用绘景：
\[
|\psi_I(t)\rangle = e^{iH_0 t/\hbar} |\psi_S(t)\rangle, \quad H'_I(t) = e^{iH_0 t/\hbar} H'(t) e^{-iH_0 t/\hbar}.
\]
演化方程：
\[
i\hbar \frac{d}{dt}|\psi_I(t)\rangle = H'_I(t) |\psi_I(t)\rangle.
\]
设 $|\psi_I(t)\rangle = \sum_n c_n(t)|n\rangle$，其中 $H_0|n\rangle = E_n|n\rangle$，则
\[
\dot{c}_f(t) = \frac{1}{i\hbar} \sum_n c_n(t) \langle f|H'_I(t)|n\rangle = \frac{1}{i\hbar} \sum_n c_n(t) \langle f|H'(t)|n\rangle e^{i\omega_{fn}t}.
\]
一级近似（设 $c_i(0)=1$，其余 $c_n(0)=0$）：
\[
c_f^{(1)}(t) = \frac{1}{i\hbar} \int_0^t \langle f|H'(t')|i\rangle e^{i\omega_{fi}t'} \, dt'.
\]
\end{derivationbox}

\subsection{四种典型时间依赖形式}

\begin{derivationbox}[四种典型时间依赖形式]
\subsubsection{常微扰：$H'(t) = V\,\Theta(t)$}

\[
c_f^{(1)}(t) = \frac{\langle f|V|i\rangle}{i\hbar} \int_0^t e^{i\omega_{fi}t'} dt' = \frac{\langle f|V|i\rangle}{i\hbar} \frac{e^{i\omega_{fi}t}-1}{i\omega_{fi}} = \frac{\langle f|V|i\rangle}{\hbar\omega_{fi}} \bigl(1 - e^{i\omega_{fi}t}\bigr).
\]
概率：
\[
P_{i\to f}(t) = \frac{4|\langle f|V|i\rangle|^2}{\hbar^2} \frac{\sin^2(\omega_{fi}t/2)}{\omega_{fi}^2}.
\]
长时间极限 ($t\to\infty$)：$\dfrac{\sin^2(\omega t/2)}{\omega^2} \to \dfrac{\pi t}{2}\delta(\omega)$，得
\[
\Gamma = \frac{2\pi}{\hbar} |\langle f|V|i\rangle|^2 \rho(E_f)
\]
即费米黄金规则，其中 $\rho(E_f)$ 是末态密度。

\subsubsection{简谐微扰：$H'(t) = V\cos(\omega t)\,\Theta(t)$}

利用 $\cos(\omega t) = \frac{1}{2}(e^{i\omega t} + e^{-i\omega t})$，积分给出两项：
\[
c_f^{(1)}(t) = -\frac{\langle f|V|i\rangle}{2\hbar} \left[ \frac{e^{i(\omega_{fi}+\omega)t}-1}{\omega_{fi}+\omega} + \frac{e^{i(\omega_{fi}-\omega)t}-1}{\omega_{fi}-\omega} \right].
\]
长时间下分母小的项主导：
\begin{itemize}
\item $\omega_{fi} \approx \omega$ ($E_f \approx E_i + \hbar\omega$)：吸收一个能量子
\item $\omega_{fi} \approx -\omega$ ($E_f \approx E_i - \hbar\omega$)：受激发射
\end{itemize}
概率与 $\dfrac{\sin^2[(\omega_{fi}\mp\omega)t/2]}{(\omega_{fi}\mp\omega)^2}$ 成正比，在共振处最大。

\subsubsection{指数开关：$H'(t) = V e^{-t/T}\,\Theta(t)$}

\[
c_f^{(1)}(\infty) = \frac{\langle f|V|i\rangle}{i\hbar} \int_0^\infty e^{-t/T} e^{i\omega_{fi}t} dt = \frac{\langle f|V|i\rangle}{i\hbar} \cdot \frac{1}{\frac{1}{T} - i\omega_{fi}}.
\]
概率：
\[
P_{i\to f} = \frac{|\langle f|V|i\rangle|^2}{\hbar^2} \cdot \frac{T^2}{1 + \omega_{fi}^2 T^2}.
\]
物理极限：
\begin{itemize}
\item $T\to\infty$（慢开关）：$P \propto T^2$，跃迁被压制（绝热定理）
\item $T\to 0$（快脉冲）：$P\to 0$，作用时间太短
\item $\omega_{fi}T \ll 1$（准静态）：$P \approx |\langle f|V|i\rangle|^2 T^2/\hbar^2$
\item $\omega_{fi}T \gg 1$（远离共振）：$P \propto 1/\omega_{fi}^2$，快速振荡平均掉
\end{itemize}

\subsubsection{$\delta$-脉冲：$H'(t) = V\,\delta(t)$}

\[
c_f^{(1)}(t) = \frac{\langle f|V|i\rangle}{i\hbar} \int_{-\epsilon}^{+\epsilon} \delta(t') e^{i\omega_{fi}t'} dt' = \frac{\langle f|V|i\rangle}{i\hbar}.
\]
概率 $P = |\langle f|V|i\rangle|^2/\hbar^2$，与时间无关。物理上对应一个瞬时踢（kick），系统状态被投影到新基上。
\end{derivationbox}

\subsection{电偶极跃迁选择定则}

\begin{derivationbox}[电偶极跃迁选择定则]
光场与原子的相互作用在电偶极近似下：$H' = -\vec{d}\cdot\vec{E}$，其中 $\vec{d} = -e\vec{r}$ 是电偶极矩算符。

\subsubsection{选择定则的代数来源}

电偶极算符 $\vec{r} = (x,y,z)$ 在角动量表象中是秩为 1 的不可约张量（向量算符）。根据 Wigner--Eckart 定理，$\langle n'l'm'|\vec{r}|nlm\rangle \neq 0$ 要求：
\begin{itemize}
\item $\Delta l = l' - l = \pm 1$（向量算符只改变角动量 1 个单位）
\item $\Delta m = m' - m = 0, \pm 1$（取决于光的偏振方向）
\item 宇称必须改变 ($\vec{r}$ 是奇宇称算符)
\end{itemize}

\textbf{考场速记：}电偶极跃迁要求 $l$ 改变 $\pm 1$，$m$ 改变 $0$ 或 $\pm 1$，且初末态宇称相反。

\subsubsection{选择定则的积分证明（以 $z$ 方向为例）}

$z = r\cos\theta \propto rY_{10}(\theta,\phi)$。矩阵元：
\[
\langle n'l'm'|z|nlm\rangle = \int R_{n'l'}(r) r R_{nl}(r) r^2 dr \int Y_{l'm'}^* Y_{10} Y_{lm} \, d\Omega.
\]
角向积分 $\int Y_{l'm'}^* Y_{10} Y_{lm} \, d\Omega$ 非零的条件是：
\begin{enumerate}
\item $m' = m + 0 = m$ ($Y_{10}$ 的磁量子数为 0)
\item 三角关系 $|l-1| \leq l' \leq l+1$，即 $l' = l \pm 1$（因为 $Y_{10}$ 的角量子数为 1）
\end{enumerate}
对 $x,y$ 方向，$x \pm iy \propto Y_{1,\pm 1}$，给出 $\Delta m = \pm 1$。
\end{derivationbox}
